\documentclass[11pt]{article}

\usepackage[margin=1in]{geometry}
\usepackage{amsmath,amssymb,amsthm}
\usepackage{graphicx}
\usepackage{float}
\usepackage[hidelinks]{hyperref}

\newtheorem{theorem}{Theorem}
\newtheorem{lemma}{Lemma}
\newtheorem{remark}{Remark}

\newcommand{\C}{\mathbb C}
\newcommand{\R}{\mathbb R}
\newcommand{\Z}{\mathbb Z}
\newcommand{\GammaSL}{\operatorname{SL}_3(\Z)}
\newcommand{\GSL}{\operatorname{SL}_3(\R)}
\newcommand{\Prob}{\operatorname{Prob}}
\newcommand{\supp}{\operatorname{supp}}

\title{A Counterexample to Removing the Invariant-Measure Hypothesis\\
in the Fejer-to-AP Theorem}
\author{\texttt{agent\_lane\_03}}
\date{2026-07-18}

\begin{document}
\maketitle

\section*{Source and Question}

The source paper is Anshu, Amrutam, and Karmakar,
\emph{Fejer property and Galois correspondence for groupoid \(C^*\)-algebras},
arXiv:2511.14474 \cite{AAK2025}.  Theorem 3.4 proves that if
\(\Gamma\ltimes X\) has the Fejer property and \(X\) admits a
\(\Gamma\)-invariant measure, then \(\Gamma\) has the Haagerup--Kraus
approximation property (AP).  The following remark then appears.

\begin{quote}
Remark 3.5.  We do not know if the assumption that \(\mu\) is invariant can be
removed in general.
\end{quote}

\begin{figure}[H]
\centering
\includegraphics[width=0.96\linewidth,height=0.56\textheight,keepaspectratio]{figures/open_problem_crop.png}
\caption{Theorem 3.4 and Remark 3.5 in arXiv:2511.14474, p. 7.}
\end{figure}

The answer is negative.

\section*{Proof Intuition}

The invariant measure in Theorem 3.4 is used to average a groupoid Fejer net
over the unit space and produce a Fejer net on the group itself.  Without such
an averaging functional, topologically amenable actions give immediate
counterexamples: an exact group can act amenably on a compact space even when
the group does not have AP.

The group \(\operatorname{SL}_3(\Z)\) is the standard exact non-AP example.  Its
action on the flag manifold \(\operatorname{SL}_3(\R)/P\), where \(P\) is a
minimal parabolic subgroup, is topologically amenable because \(P\) is amenable.
Adding a profinite left-translation factor makes the action free while
preserving amenability.  The resulting transformation groupoid has bounded
Fejer property, but the acting group still lacks AP.

\section*{The Counterexample}

\begin{theorem}
There is a compact metrizable \(\GammaSL\)-space \(X\) such that the
transformation groupoid \(\GammaSL\ltimes X\) is second countable, Hausdorff,
etale, principal, and has the bounded Fejer property, while \(\GammaSL\) does
not have AP.

Consequently, the invariant-measure hypothesis in Theorem 3.4 of
\cite{AAK2025} cannot be removed.
\end{theorem}

\begin{proof}
Let \(\Gamma=\GammaSL\) and let \(G=\GSL\).  Let \(P<G\) be a minimal parabolic
subgroup; concretely, one may take the upper triangular subgroup of \(G\).
The quotient
\[
   Y=G/P
\]
is the compact metrizable full flag manifold.

The subgroup \(P\) is amenable.  Hence the action \(G\curvearrowright G/P\) is
topologically amenable, and restricting the action to the discrete subgroup
\(\Gamma<G\) preserves topological amenability.

Let \(Z=\widehat{\Gamma}\) be the profinite completion of \(\Gamma\), with
\(\Gamma\) acting by left translation.  Since \(\Gamma\) is finitely generated
and residually finite, \(Z\) is compact metrizable and the natural map
\(\Gamma\to Z\) is injective.  The left-translation action on \(Z\) is free:
if \(g z=z\) for some \(z\in Z\), then \(g=e\) in the group \(Z\), hence
\(g=e\) in \(\Gamma\).

Set
\[
   X=Y\times Z
\]
with the diagonal \(\Gamma\)-action.  The action remains topologically amenable:
if \(m_i:Y\to\Prob(\Gamma)\) are continuous approximately equivariant maps
witnessing amenability of \(\Gamma\curvearrowright Y\), then
\[
   (y,z)\longmapsto m_i(y)
\]
witnesses amenability of \(\Gamma\curvearrowright X\).  It is also free because
the second coordinate action is free.  Therefore \(\Gamma\ltimes X\) is a
second countable Hausdorff etale principal transformation groupoid.

It remains only to record that this groupoid has the bounded Fejer property.
For a topologically amenable action of a discrete group on a compact space,
there is a net \(h_i\in C_c(\Gamma\ltimes X)\) of positive-type functions such
that \(h_i\to 1\) uniformly on compact subsets of \(\Gamma\ltimes X\).  The
pointwise multipliers \(M_{h_i}\) are completely positive contractions on
\[
   C_r^*(\Gamma\ltimes X)\cong C(X)\rtimes_r\Gamma.
\]
For an element \(a=\sum_{s\in F} a_s u_s\in C_c(\Gamma\ltimes X)\), where
\(F\subset\Gamma\) is finite, uniform convergence on \(F\times X\) gives
\[
\begin{aligned}
\|M_{h_i}(a)-a\|
&\leq \sum_{s\in F}\|(h_i(s,\cdot)-1)a_s\|_{C(X)}  \\
&\leq \sum_{s\in F}\|a_s\|_{C(X)}
   \sup_{(s,x)\in F\times X}|h_i(s,x)-1|
   \longrightarrow 0.
\end{aligned}
\]
Since \(C_c(\Gamma\ltimes X)\) is dense in \(C_r^*(\Gamma\ltimes X)\) and the
multipliers have cb-norm at most \(1\), the convergence holds for every element
of the reduced crossed product.  Thus \(\Gamma\ltimes X\) has the bounded Fejer
property.

On the other hand, \(\operatorname{SL}_3(\Z)\) does not have AP, by the
Lafforgue--de la Salle theorem as cited in \cite{AAK2025} and in the standard
literature \cite{HdL2013}.  Therefore the implication
\[
   \Gamma\ltimes X \text{ has Fejer property}
   \quad\Longrightarrow\quad
   \Gamma \text{ has AP}
\]
is false without an invariant-measure hypothesis.  This is precisely the
proposed removal in Remark 3.5.
\end{proof}

\begin{remark}
The example necessarily has no \(\Gamma\)-invariant probability measure on
\(X\).  Indeed, if such a measure existed, Theorem 3.4 of the source paper
would apply to the Fejer groupoid \(\Gamma\ltimes X\) and would force
\(\GammaSL\) to have AP, contradicting the final paragraph of the proof.
\end{remark}

\section*{Verification Notes}

No numerical computation is used.  The proof depends on the following standard
inputs:
\begin{enumerate}
\item \(\operatorname{SL}_3(\Z)\) is finitely generated, residually finite, and
does not have AP.
\item If \(P\) is a minimal parabolic subgroup of \(\operatorname{SL}_3(\R)\),
then \(P\) is amenable and the action on \(G/P\) is topologically amenable.
\item A topologically amenable action of a discrete group on a compact space
gives compactly supported positive-type Herz--Schur multipliers converging to
\(1\), hence the bounded Fejer property for the transformation groupoid.
\end{enumerate}

\section*{Scope and Novelty Check}

This packet answers only Remark 3.5/Theorem 3.4 of the source paper.  It does
not address Question 1 of Fuller--Karmakar about Fourier recovery for arbitrary
groupoids, and it does not challenge the source paper's Galois correspondence
or bimodule theorems.

The run's cheap indexes were searched for arXiv:2511.14474, the source title,
``Fejer'', ``invariant measure'', ``removed in general'', ``amenable action'',
``transformation groupoid'', ``SL3'', and ``AP''.  No prior packet for this
paper or route was found.  A bounded web search on 2026-07-18 for exact and
nearby phrases including ``Fejer property invariant measure removed groupoid'',
``Gamma rtimes X Fejer property invariant measure AP'', and
``SL3(Z) approximation property exact amenable action'' found the source paper
and standard background on exact non-AP groups and amenable actions, but no
existing answer to Remark 3.5.

\section*{Human Review Recommendation}

Recommended verdict: likely valid counterexample.  Reviewers should mainly
check the standard amenability-to-positive-type-multiplier passage for the
specific Fejer definition in the source paper.  The obstruction to dropping
the invariant-measure hypothesis is otherwise direct: \(\Gamma\ltimes X\) has
bounded Fejer property while \(\Gamma=\operatorname{SL}_3(\Z)\) lacks AP.

\begin{thebibliography}{9}

\bibitem{AAK2025}
Anshu, Tattwamasi Amrutam, and Pradyut Karmakar,
\emph{Fejer property and Galois correspondence for groupoid \(C^*\)-algebras},
arXiv:2511.14474, 2025.

\bibitem{ADR2000}
Claire Anantharaman-Delaroche and Jean Renault,
\emph{Amenable groupoids},
Monographies de L'Enseignement Mathematique 36, 2000.

\bibitem{HdL2013}
Uffe Haagerup and Tim de Laat,
\emph{Simple Lie groups without the approximation property},
Duke Mathematical Journal 162 (2013), no. 5, 925--964.

\bibitem{BrownOzawa2008}
Nathanial P. Brown and Narutaka Ozawa,
\emph{\(C^*\)-algebras and finite-dimensional approximations},
Graduate Studies in Mathematics 88, American Mathematical Society, 2008.

\end{thebibliography}

\end{document}
